% 第 9 题三视图：按固定源 PDF 矢量路径换算（1 pt → 0.03515 cm）。
% 原点取正视图左下角（源页 (531.30, 382.71) pt），y 向上。网格 25.5 pt = 2 cm。
% 正视图：L 形（左柱 2×2 在下、右柱 2×4 通高；中竖贯通、中横仅左半——按源矢量路径复核）；
% 侧视图：与正视图互为镜像（左柱通高、右柱仅下段），右柱右棱下半为虚线（源内线型如此，照录）；
% 俯视图：左列满深 2×4、右上 2×2 正方形。
% 尺寸标注：横向 4 条（正视图两格 + 侧视图两格）@y=56.10、竖向 2 条 @x=56.10（均两段让位给数字）；
% 6 条竖直刻度线、3 条水平刻度线；12 个实心箭头（一律朝外）；「2」「正视图」「侧视图」「俯视图」共 9 个标签（照源字形框中心）。
\begin{tikzpicture}[line cap=round,line join=round,font=\small,x=0.03515cm,y=0.03515cm]
% ===== 正视图（左格）=====
\draw (0,0)--(51,0)--(51,51)--(25.5,51)--(25.5,25.5)--(0,25.5)--cycle;
\draw (25.5,0)--(25.5,25.5);                 % 中竖下半（实线）
% ===== 侧视图（右格，右移 76.5）=====
\draw (76.5,0)--(127.51,0)--(127.51,25.5)--(102,25.5)--(102,51)--(76.5,51)--cycle;
\draw (76.5,25.5)--(102,25.5);               % 中横左半（实线）
\draw[dash pattern=on 2.99pt off 2.99pt] (102,0)--(102,25.5);  % 中竖下半（虚线，照源）
% ===== 俯视图 =====
\draw (0,-25.5)--(25.5,-25.5)--(25.5,-76.51)--(0,-76.51)--cycle;   % 左列满深
\draw (25.5,-25.5)--(51,-25.5)--(51,-51.01)--(25.5,-51.01);        % 右上正方形（第四边与左列右边重合）
% ===== 尺寸线（横向 4 条 @y=56.10，两段让位）=====
\draw (3.58,56.10)--(21.91,56.10);   \draw (29.08,56.10)--(47.41,56.10);       % 正视图两格
\draw (80.09,56.10)--(98.42,56.10);  \draw (105.59,56.10)--(123.92,56.10);     % 侧视图两格
% ===== 尺寸线（竖向 2 条 @x=56.10，两段让位）=====
\draw (56.10,29.09)--(56.10,47.42);  % 上段（高 2）
\draw (56.10,3.58)--(56.10,21.91);   % 下段（高 2）
% ===== 竖直刻度线（6 条，@y 53.55–61.20）=====
\draw (0,53.55)--(0,61.20);         \draw (25.5,53.55)--(25.5,61.20);
\draw (51,53.55)--(51,61.20);       \draw (76.5,53.55)--(76.5,61.20);
\draw (102,53.55)--(102,61.20);     \draw (127.51,53.55)--(127.51,61.20);
% ===== 水平刻度线（3 条，@x 53.55–61.20）=====
\draw (53.55,0)--(61.20,0);
\draw (53.55,25.5)--(61.20,25.5);
\draw (53.55,51)--(61.20,51);
% ===== 实心箭头（12 个，底 2.98 长 3.98，一律朝外）=====
\fill (0,56.10)--(3.98,54.61)--(3.98,57.59)--cycle;             % 正左格：左端←
\fill (25.5,56.10)--(21.52,54.61)--(21.52,57.59)--cycle;       % 正左格：右端→
\fill (25.5,56.10)--(29.48,54.61)--(29.48,57.59)--cycle;       % 正右格：左端←
\fill (51,56.10)--(47.02,54.61)--(47.02,57.59)--cycle;         % 正右格：右端→
\fill (76.5,56.10)--(80.48,54.61)--(80.48,57.59)--cycle;       % 侧左格：左端←
\fill (102,56.10)--(98.02,54.61)--(98.02,57.59)--cycle;        % 侧左格：右端→
\fill (102,56.10)--(105.98,54.61)--(105.98,57.59)--cycle;      % 侧右格：左端←
\fill (127.51,56.10)--(123.53,54.61)--(123.53,57.59)--cycle;   % 侧右格：右端→
\fill (56.10,51)--(54.61,47.02)--(57.59,47.02)--cycle;         % 上段：上端↑
\fill (56.10,25.5)--(54.61,29.48)--(57.59,29.48)--cycle;       % 上段：下端↓
\fill (56.10,25.5)--(54.61,21.52)--(57.59,21.52)--cycle;       % 下段：上端↑
\fill (56.10,0)--(54.61,3.98)--(57.59,3.98)--cycle;            % 下段：下端↓
% ===== 标签（照源字形框中心，显式坐标）=====
\node[anchor=center,inner sep=0] at (12.96,60.43) {\(2\)};
\node[anchor=center,inner sep=0] at (38.47,60.43) {\(2\)};
\node[anchor=center,inner sep=0] at (89.49,60.43) {\(2\)};
\node[anchor=center,inner sep=0] at (115.01,60.43) {\(2\)};
\node[anchor=center,inner sep=0] at (61.43,37.46) {\(2\)};
\node[anchor=center,inner sep=0] at (61.43,11.95) {\(2\)};
\node[anchor=center,inner sep=0] at (25.71,-12.21) {正视图};
\node[anchor=center,inner sep=0] at (102.24,-12.19) {侧视图};
\node[anchor=center,inner sep=0] at (25.71,-88.67) {俯视图};
\end{tikzpicture}
